%%%%%%%%%%%%%%%%%%%%%%%%%%%%%%%%%%%%%%%%
% RESUMO %% obrigat�rio
\begin{resumo}
%% neste arquivo resumo.tex o texto do resumo e as palavras-chave t�m que ser em Portugu�s para os documentos escritos em Portugu�s
%% o texto do resumo e as palavras-chave t�m que ser em Ingl�s para os documentos escritos em Ingl�s
%% os nomes dos comandos \begin{resumo}, \end{resumo}, \palavraschave e \palavrachave n�o devem ser alterados
\hypertarget{estilo:resumo}{} %% uso para este Guia


O NanoSatC-Br1 � um microssat�lite brasileiro, ele � um CubeSat executado no �mbito do Conv�nio do Instituto Nacional de Pesquisas Espaciais (INPE), atrav�s de sua subunidade o Centro Regional Sul de Pesquisas Espaciais (CRS) com a Universidade Federal de Santa Maria (UFSM). O OSIRIS-REx � uma miss�o de ci�ncia planet�ria, que consiste em estudar e coletar amostras do asteroide 101955 Bennu, um asteroide carbon�ceo, chegando a Terra de volta em 2023. O objetivo deste trabalho � mostrar, para cada miss�o, o objetivo, diagramas, fotos, funcionamento, e  equipamentos utilizados.




\palavraschave{
	\palavrachave{Miss�es}
}


\end{resumo}